\section{Future works}\label{future-works}

Sebbene il progetto raggiunga gli obiettivi didattici prefissati in termini di replica distribuita, convergenza eventuale e osservabilit\`a, l'architettura realizzata lascia spazio a diverse estensioni significative, sia dal punto di vista ingegneristico sia da quello sperimentale.

\subsection{Sicurezza e hardening operativo}

La prima area di evoluzione riguarda la sicurezza. Nella versione attuale il sistema opera in un contesto controllato e non implementa autenticazione dei peer, protezione dell'API HTTP n\'e cifratura dei canali. Un'evoluzione naturale consisterebbe nell'introdurre:
\begin{itemize}
  \item autenticazione reciproca tra nodi;
  \item protezione degli endpoint HTTP mediante token o mTLS;
  \item validazione pi\`u stringente dell'identit\`a logica del mittente rispetto al canale di trasporto;
  \item policy di autorizzazione per distinguere ruoli osservativi e amministrativi.
\end{itemize}

Queste modifiche renderebbero il sistema pi\`u vicino a uno scenario reale, nel quale la tolleranza ai guasti deve convivere con requisiti minimi di trust e protezione del piano di controllo.

\subsection{Persistenza e storico dei dati}

Attualmente i nodi mantengono soltanto stato volatile in memoria. Questa scelta \`e coerente con il focus del progetto, ma limita la capacit\`a del sistema di sopravvivere a riavvii estesi o di offrire analisi temporali di lungo periodo. Un'estensione rilevante sarebbe quindi l'introduzione di:
\begin{itemize}
  \item persistenza locale dei record sensore e dei delta recenti;
  \item snapshot periodiche del metadata di membership e topology;
  \item meccanismi di restore all'avvio per ridurre il costo di bootstrap;
  \item archiviazione storica separata per analisi offline e confronto tra run sperimentali.
\end{itemize}

Questa evoluzione aprirebbe anche la possibilit\`a di confrontare in modo pi\`u rigoroso comportamento \emph{cold start}, tempi di recovery e perdita temporanea di informazione dopo crash estesi.

\subsection{Semantiche di replica pi\`u ricche}

Il criterio Last-Writer-Wins adottato \`e semplice, deterministico e adatto al contesto del progetto, ma costituisce anche una semplificazione. In prospettiva sarebbe interessante valutare:
\begin{itemize}
  \item l'uso di CRDT specifici per particolari tipi di dato;
  \item merge policies differenziate per categoria di sensore;
  \item version vector o metadata causali per distinguere meglio concorrenza reale e semplice arrivo tardivo;
  \item strategie adattive tra replica incrementale e full sync in funzione del grado di divergenza osservato.
\end{itemize}

Queste estensioni consentirebbero di studiare in modo pi\`u approfondito il compromesso fra semplicit\`a implementativa, quantit\`a di metadata scambiati e qualit\`a semantica della convergenza.

\subsection{Sperimentazione su topologie pi\`u ampie}

Un'altra direzione naturale riguarda l'ampliamento della scala sperimentale. Le topologie attuali permettono gi\`a di osservare fenomeni interessanti, ma un'estensione a cluster pi\`u grandi o con connettivit\`a pi\`u eterogenea renderebbe possibile:
\begin{itemize}
  \item valutare l'impatto della densit\`a della topologia sul traffico di gossip;
  \item studiare la sensibilit\`a del failure detector in reti meno regolari;
  \item confrontare fan-out, intervalli di sync e dimensione dei delta buffer su carichi superiori;
  \item verificare la leggibilit\`a e la sostenibilit\`a dell'osservabilit\`a al crescere del numero di nodi.
\end{itemize}

Da un punto di vista sperimentale, sarebbe inoltre utile affiancare metriche quantitative sistematiche su latenze di convergenza, volume di traffico e costi di recovery.

\subsection{Osservabilit\`a avanzata e strumenti di analisi}

Il progetto dispone gi\`a di una buona base di introspection read-only, ma essa potrebbe essere ulteriormente estesa tramite:
\begin{itemize}
  \item esportazione di metriche verso stack standard di monitoring;
  \item alerting automatico su soglie critiche di stale data o \texttt{DELTA\_UNAVAILABLE};
  \item tracciamento temporale persistente degli eventi di control plane;
  \item comparazione automatizzata tra differenti configurazioni di replica e failure detection.
\end{itemize}

In questo modo il sistema diventerebbe non solo un prototipo funzionale, ma anche una piattaforma sperimentale pi\`u matura per analizzare in maniera sistematica le conseguenze delle scelte architetturali.

\clearpage
